% Ad Quintum Fratrem 2.7 (ad-quintum-fratrem-02-07) --- English translation
% Translated by Claude (Anthropic) from the Latin (Perseus).
% Style guide: STYLE.md.

\ciceroLetterOpener{MARCVS QVINTO FRATRI salutem}

\ciceroSection{1} I had supposed my book would please you; that it pleased so much as you write delights me very much. As for your reminding me of our Urania and urging me to remember the speech of Jupiter that stands at the end of that book --- yes, I do remember; and I wrote all that more for myself than for the rest.

\ciceroSection{2} But here is the news. The day after you set out, I came late at night with Vibullius to Pompey; and when I had spoken to him about those works and inscriptions of yours, he answered me very kindly and gave me great hopes. He said that he wished to speak with Crassus and urged me to do the same. As consul I escorted Crassus from the Senate to his house. He took the matter in hand, and said that there was something Clodius wished at this time to obtain through him and through Pompey: he supposed, if I did not stand in the man's way, that I could obtain what I wanted without a fight. The whole business I committed to him and said I would be in his hands. The young Publius Crassus, our most devoted as you know, was party to this conversation. What Clodius wants is some legation \textit{(if not by Senate, then by people)} either ``free,'' or to Byzantium, or to Brogitarus, or both --- a thing full of money. About which I am not greatly exercised, even if I get less than I want. Pompey has spoken with Crassus all the same. They seem to have taken the business in hand. If they pull it off, excellent; if not, let us go back to our own Jupiter.

\ciceroSection{3} On the third day before the Ides of February the Senate passed its decree on bribery in the formula proposed by Afranius, the formula I had supported when you were present. But to the deep groans of the Senate, the consuls did not push through the views of those who, while assenting to Afranius, had added the rider that the praetors should be elected so as to have sixty days as private citizens. On that day they rejected Cato outright. In short --- they hold everything, and they want everyone to know it.
